\subsection{Quantifiers, negation, and the standard vocabulary}\label{sub:quant}
% Defined once, here, at the top of the sheet; every later section \ref's back to this
% box instead of re-explaining "compact" or "dense" in passing.
\D{the vocabulary every exercise uses --- defined here, used everywhere}{
\begin{bul}
\item \hd{Intervals:} \emph{open} $(a,b)$ (no endpoints), \emph{closed} $[a,b]$ (both), half-open $[a,b)$. \hd{$|I|$} $=$ length.
\item \hd{Bounded:} $\exists C:|x|\le C$ for all $x$ in the set. Bounded above / below separately.
\item \hd{Compact} \de{kompakt} in $\Rl$ $\iff$ \textbf{closed and bounded} (Heine--Borel). This one word is the hypothesis of IVT, Weierstrass and Heine (\S\ref{sub:bigthree}); drop it and all three fail.
\item \hd{Neighbourhood / $\eps$-ball} of $x$: the interval $(x-\eps,x+\eps)$. \hd{interior point:} some $\eps$-ball fits entirely inside the set; \hd{boundary point:} every $\eps$-ball meets the set \emph{and} its complement.
\item \hd{Null sequence:} $a_n\to0$. \hd{Subsequence} $(a_{n_k})$: indices \textbf{strictly increasing}, $n_1<n_2<\cdots$ (so $n_k\ge k$).
\item \hd{Dense} in $\Rl$: every interval, however short, contains a point of the set. $\Q$ and $\Rl\setminus\Q$ both are.
\end{bul}
}
\R{negate any definition mechanically}{
\hd{Rule:} swap $\forall\leftrightarrow\exists$ left to right, \textbf{keep the order}, negate the final inequality ($<$ becomes $\ge$). \warn The \hd{order} is the whole content: $\forall x\,\exists\delta$ (continuity, $\delta$ may depend on $x$) is a \emph{different theorem} from $\exists\delta\,\forall x$ (uniform, \textbf{one} $\delta$ for all $x$).
\begin{bul}
\item \hd{$a_n\not\to a$:} $\exists\eps>0\ \forall N\ \exists n\ge N:\ |a_n-a|\ge\eps$.
\item \hd{not Cauchy:} $\exists\eps>0\ \forall N\ \exists m,n\ge N:\ |a_n-a_m|\ge\eps$.
\item \hd{$f$ not continuous at $x_0$:} $\exists\eps>0\ \forall\delta>0\ \exists x:|x-x_0|<\delta$ \emph{and} $|f(x)-f(x_0)|\ge\eps$.
\item \hd{$f$ not uniformly continuous:} $\exists\eps>0\ \forall\delta>0\ \exists x,y:|x-y|<\delta$ and $|f(x)-f(y)|\ge\eps$.
\item \hd{$f_n\not\rightrightarrows f$:} $\exists\eps>0\ \forall N\ \exists n\ge N\ \exists x:|f_n(x)-f(x)|\ge\eps$.
\end{bul}
\hd{In practice} the last two become: build $x_n,y_n$ resp.\ $x_n$ and show the gap does not vanish.
}

\subsection{sup / inf / max / min \de{Supremum, Infimum}}\label{sub:supinf}
\D{the four words}{
$M\subseteq\Rl$ nonempty. \hd{upper bound} $s$: $x\le s\ \forall x\in M$.
\hd{$\sup M$} = least upper bound. \hd{$\max M$} = $\sup M$ \emph{only if} $\sup M\in M$. Mirrored for $\inf/\min$.\\
\hd{Completeness axiom:} every nonempty $M$ bounded above has $\sup M\in\Rl$. This is the one property $\Q$ lacks and the engine behind IVT, Bolzano--Weierstrass and monotone convergence.
}
\F{$\eps$-characterisation (use it in every sup proof)}{
$s=\sup M\iff$ (i) $x\le s\ \forall x\in M$, \ and (ii) $\forall\eps>0\ \exists x\in M:\ x>s-\eps$.\\
For a function: $\sup_{x\in D}f(x)$ is the sup of the image set $f(D)$; it is a $\max$ iff attained.
}
\R{find $\sup/\inf$ and prove it}{
\begin{rcp}
\item Guess $s$ (the largest/smallest value the formula can approach).
\item \hd{Bound:} show $x\le s$ for all $x\in M$ by algebra.
\item \hd{Least:} given $\eps>0$ produce $x\in M$ with $x>s-\eps$ (usually choose $n>1/\eps$).
\item Attained? $s\in M$ \ra\ $\max=s$; otherwise $\max$ does not exist.
\end{rcp}
}
% D-INFK FS25 1.MC2 and 2.A1 were both exactly this, and neither was answerable from
% the sup/inf recipe alone: they ask about a SET, not a sequence.
\R{a \emph{set} is given by a formula: open? closed? bounded? $\sup$? $\inf$?}{
\hd{Open:} every point has an $\eps$-ball still inside $M$. \hd{Closed:} the complement is open $\iff$ every convergent sequence \emph{in} $M$ has its limit in $M$.
\begin{rcp}
\item \hd{Given by an inequality on a continuous $g$:} $\{g>c\}$ and $\{g<c\}$ are \textbf{open}; $\{g\ge c\}$, $\{g\le c\}$, $\{g=c\}$ are \textbf{closed}. Solve the inequality into intervals and read it off.
\item \hd{Given by a sequence formula} $\{a_n:n\in\N\}$: check monotonicity. Monotone decreasing \ra\ $\sup=\max=a_1$ (attained), $\inf=\lim a_n$ (\textbf{not} attained). Increasing \ra\ mirrored. Not monotone \ra\ split by parity first (\S\ref{sub:accum}).
\item \hd{Bounded?} both $\sup$ and $\inf$ finite. Unbounded above \ra\ write $\sup=+\infty$.
\end{rcp}
\warn ``\emph{neither open nor closed}'' is the usual trap answer: $[a,b)$ qualifies, but $\Rl$ and $\emptyset$ are \textbf{both} open and closed, and a finite set is closed.\\
\hd{Ex.} $A=\{x:|x+1|^3>2\}=\{|x+1|>\sqrt[3]2\}=(-\infty,-1-\sqrt[3]2)\cup(-1+\sqrt[3]2,\infty)$: \textbf{open}, unbounded, $\sup A=+\infty$, $\inf A=-\infty$.\\
\hd{Ex.} $B=\left\{\frac1{n+1}-\frac1{n+2}\right\}=\left\{\frac1{(n+1)(n+2)}\right\}$: decreasing \ra\ $\max$ at the first index, $\inf=0$ not attained.
}
\F{Archimedes and density}{
\hd{Archimedean:} $\forall x\in\Rl\ \exists n\in\N:n>x$; equivalently $\frac1n\to0$. $\N$ is unbounded in $\Rl$.\\
\hd{$\Q$ dense:} between any $a<b$ there is a rational \emph{and} an irrational. Used to build the Dirichlet counterexamples.
}
\X{sup/inf in multiple choice}{
\begin{bul}
\item $\sup$ exists as a \emph{real number} only for a bounded-above set. Unbounded \ra\ write $\sup=+\infty$, i.e.\ no real sup.
\item $\min M$ need \textbf{not} be an accumulation point of the sequence: $x_0=0$, $x_n=1+\frac1n$ attains its min once.
\item \claimT{$\inf M=\sup M$ \ra\ $M$ is a single point \ra\ the sequence is constant \ra\ converges.}
\item monotone \emph{and bounded} \ra\ converges to $\sup$ (incr.) / $\inf$ (decr.). Monotone alone may run to $\pm\infty$.
\end{bul}
}

\subsection{Absolute value and inequalities}\label{sub:ineq}
\R{solve $|\cdot|$ inequalities}{
$|x-a|<r\iff x\in(a-r,a+r)$;\quad $|x-a|>r\iff x<a-r$ or $x>a+r$.\\
$|f(x)|<g(x)\iff -g(x)<f(x)<g(x)$ (and $g>0$).\\
General: split at every zero of the inside, solve on each piece, intersect with that piece.
}
\F{the inequalities you actually use}{
\begin{bul}
\item \hd{Triangle:} $|a+b|\le|a|+|b|$; \hd{reverse:} $\big||a|-|b|\big|\le|a-b|$.
\item \hd{Bernoulli:} $(1+x)^n\ge1+nx$ for $x\ge-1$.
\item \hd{AM--GM:} $\sqrt{ab}\le\frac{a+b}2$ $(a,b\ge0)$.
\item \hd{Young:} $ab\le\frac{a^2}2+\frac{b^2}2$, sharper $ab\le\frac{\eps a^2}{2}+\frac{b^2}{2\eps}$.
\item \hd{Cauchy--Schwarz:} $|\langle x,y\rangle|\le\|x\|\|y\|$.
\end{bul}
\hd{Always available:} $|\sin t|\le1$, $|\sin t|\le|t|$ (strict for $t\ne0$), $|\cos t|\le1$, $e^t\ge1+t$, $\ln t\le t-1$.
}
\R{estimate without panic}{
Target is always $|\text{expr}|\le(\text{something}\to0)$.
\begin{bul}
\item Numerator: make it \emph{bigger} (drop negative terms, bound $\sin/\cos$ by 1).
\item Denominator: make it \emph{smaller} (drop positive terms). Never the reverse.
\item Roots: conjugate, $\sqrt A-\sqrt B=\frac{A-B}{\sqrt A+\sqrt B}$.
\item \hd{Sums:} bound every term by the largest, then multiply by the count.
\end{bul}
}
\F{small-$x$ behaviour: the six approximations, and the inequalities that prove them}{
For $x\to0$ \de{and only then}, each function may be replaced by its first Taylor term; the error is one order higher (full list \S\ref{sub:seriesops}, \S\ref{sub:taylorlist}).
\begin{tsTabular}{@{}l@{\ \ }l@{\quad}l@{}}
\hd{$\sin x\approx x$} & $x-\frac{x^3}6+O(x^5)$ & \hd{exact:} $|\sin x|\le|x|$ always\\
\hd{$\tan x\approx x$} & $x+\frac{x^3}3+O(x^5)$ & $\tan x\ge x$ on $[0,\frac\pi2)$\\
\hd{$1-\cos x\approx\frac{x^2}2$} & $\frac{x^2}2-\frac{x^4}{24}+O(x^6)$ & $0\le1-\cos x\le\frac{x^2}2$\\
\hd{$e^x\approx1+x$} & $1+x+\frac{x^2}2+O(x^3)$ & \hd{exact:} $e^x\ge1+x$ for \emph{all} $x$\\
\hd{$\ln(1+x)\approx x$} & $x-\frac{x^2}2+O(x^3)$ & $\frac x{1+x}\le\ln(1+x)\le x$\\
\hd{$(1+x)^\alpha\approx1+\alpha x$} & $1+\alpha x+O(x^2)$ & \hd{Bernoulli:} $\ge$ for $\alpha\ge1$
\end{tsTabular}
\warn Use the \hd{exact} column when you must \emph{prove} an estimate; the $\approx$ column only inside a limit. Substituting $x=\frac1n$ turns every row into a sequence estimate, e.g. $n\sin\frac1n\to1$.
}

\subsection{Induction \de{vollst\"andige Induktion}}\label{sub:induction}
\R{proof by induction}{
\begin{rcp}
\item \hd{Base:} verify $n=n_0$ explicitly.
\item \hd{Hypothesis:} assume $A(n)$ for one fixed $n\ge n_0$.
\item \hd{Step:} write the $n{+}1$ statement, isolate the $n$ part, substitute the hypothesis, then estimate to the target. Say where you used it.
\end{rcp}
\hd{Helpers:} $\sum_{k=1}^nk=\frac{n(n+1)}2$, $\sum_{k=1}^nk^2=\frac{n(n+1)(2n+1)}6$, $\sum_{k=0}^nq^k=\frac{1-q^{n+1}}{1-q}$, $(1+x)^n\ge1+nx$.
}

\subsection{Vectors, inner product, norms in $\Rl^n$}\label{sub:vectors}
\D{inner product, norm, distance}{
$\langle x,y\rangle=x\cdot y=\sum_{i=1}^nx_iy_i$;\quad
$\|x\|=\sqrt{\langle x,x\rangle}=\sqrt{\sum x_i^2}$;\quad
$d(x,y)=\|x-y\|=\sqrt{\sum(x_i-y_i)^2}$.
}
\F{Cauchy--Schwarz and the norm axioms}{
$|\langle x,y\rangle|\le\|x\|\,\|y\|$, with \textbf{equality iff $x,y$ are linearly dependent}. This is exactly what makes the triangle inequality true.\\
A \hd{norm} satisfies: $\|x\|\ge0$ and $\|x\|=0\iff x=0$; $\|cx\|=|c|\|x\|$; $\|x+y\|\le\|x\|+\|y\|$.\\
Hence $\|x-z\|\le\|x-y\|+\|y-z\|$ for all $x,y,z$.
}
